\section{Teoría de la Probabilidad}

\subsection{Cadenas de Markov}

Es un proceso estocástico $\{X_1, X_2 \cdots X_n\}$ en tiempo discreto. Este cumple con la propiedad de Markov:

\begin{align*}
    &P(X_{n+1} = e_j \mid X_n = e_i, \cdots X_0 = e_{0}) \\
    &= P(X_{n+1} = e_j \mid X_n = e_i)
\end{align*}

Donde el espacio de estados es $E = \{e_1, e_2\cdots e_m\}$ finito o numerable.

Aquí, la probabilidad de pasar de un estado $e_i$ a un estado $e_j$ en un solo paso está dado por:

\begin{align*}
    p_{ij} = P(X_{n+1} = e_j \mid X_n = e_i)
\end{align*}

Estas son las probabilidades de transición, usualmente organizadas en la matriz de transición $P$:

\begin{align*}
    P = \begin{bmatrix}
        p_{11} & p_{12} & \cdots & p_{1m}\\
        p_{21} & p_{22} & \cdots & p_{2m}\\
        \vdots & \vdots & \ddots & \vdots\\
        p_{m1} & p_{m2} & \cdots & p_{mm}
    \end{bmatrix}
\end{align*}

Donde las filas son el estado actual, y las columnas el estado siguiente. La suma de los elementos de cada fila deben sumar $1$.

Para $n$ pasos, se utiliza la matriz $P^{(n)}$ y se encuentra el respectivo elemento $p^{(n)}_{ij}$. La matriz $P^{(n)}$ se calcula multiplicando $n$ veces la matriz $P$ por sí misma.

También tenemos la distribución de probabilidad del sistema $\pi$. La distribución de probabiliadad del sistema tras $n$ pasos está dada por:

\begin{align*}
    \pi^{(n)} = (\pi^{(n)}_1,\ldots,\pi^{(n)}_m)
\end{align*}

Donde 

\begin{align*}
    \pi^{(n)}_j = P(X_n = e_j)
\end{align*}

Es la probabilidad de que el sistema se encuentre en el estado $e_j$ en el instante $n$ o después de $n$ pasos.

Note que esta va evolucionando, donde cada paso modifica las probabilidades con

\begin{align*}
    \pi^{(n+1)} = \pi^{(n)} P
\end{align*}

Y de forma más general se cumple que:

\begin{align*}
    \pi^{(n)} = \pi^{(0)} P^{(n)}
\end{align*}

Siendo $\pi^{(0)}$ la distribución de probabilidades inicial.